% SSSUP-FONTE-STORICA: fonti-locali/originali/2014-2020.zip
% Anno 2019, esercizio 2. Trasposizione dal sorgente LaTeX storico.
% Verificare editorialmente la trasposizione prima di considerarla definitiva.
\renewcommand{\Anno}{2019}
\renewcommand{\NumeroProblema}{2}

\renewcommand{\Testo}{
Quattro allievi del Sant'Anna (Antonio, Beatrice, Cinzia e Dario, indicati con A, B, C e D per comodità) fanno il seguente gioco, dove un intero $N\ge1$ è fissato a priori.
\begin{itemize}
\item Il gioco si svolge a squadre: A e C contro B e D
\item A inizia dicendo un numero intero da $1$ a $n_A$;
\item B continua dicendo a sua volta un numero intero da $1$ a $n_B$ e calcola la somma dei numeri detti finora;
\item poi è il turno degli altri giocatori C e D che dicono a loro volta un numero intero compreso fra $1$ e $n_C$ e fra $1$ e $n_D$, rispettivamente.
\item Poi si continua di nuovo con A e così via.
\item Vince la squadra che dice il numero che permette alla somma dei numeri detti dall'inizio del gioco di arrivare a $N$.
\end{itemize}

\begin{enumerate}
\item[(a)] Assumendo $n_A=n_B=n_C=n_D=n$ e $n\ge2$, esiste una strategia con cui una delle due squadre riesce a vincere con certezza? Se sì, quale squadra e con quale strategia? Se no, perché non esiste una strategia?
\item[(b)] Assumendo $n_A=15$, $n_B=10$, $n_C=5$, $n_D=10$ e $N=500$, esiste una strategia con cui una delle due squadre riesce a vincere con certezza? Se sì, quale squadra e con quale strategia? Se no, perché non esiste una strategia?
\end{enumerate}
}

\renewcommand{\Risposta}{}
\renewcommand{\Aiutino}{}
\renewcommand{\SoluzioneUfficiale}{}
\renewcommand{\SoluzioneRedazionale}{}
\renewcommand{\SoluzioniAlternative}{}
\renewcommand{\Osservazioni}{}
\renewcommand{\Commenti}{}
